% !TeX root = ../../Book3.tex
\section{同伦形式的 Cauchy 定理}
\label{b3:sec:1002}

局部原函数使积分在路径小变形下不变。把局部相消沿同伦方格拼接，可以对一般区域的同伦路径证明积分不变；单连通性只是全部闭路都能收缩时的一个特例。

\subsection{同伦曲线}
\label{b3:subsec:100201}

\begin{theorem}\label{b3:thm:10-homotopy-integral}
设 \(f\) 在区域 \(U\) 全纯。端点固定同伦的两条分段光滑路径积分相同；自由同伦的两条闭路积分也相同。
\end{theorem}
\begin{proof}
同伦像紧，取有限个有局部原函数的圆盘覆盖它，并用 Lebesgue 数把参数方形分成足够小的格。每个小格像包含在一个圆盘内；即使内部边只是连续路径，也可用原函数端点差定义积分。每格四边端点差之和为0。把全部小格相加，内部边方向相反而抵消，只剩方形外边。端点固定情形，两条竖边常值，积分为0；闭路自由同伦情形，两条竖边是同一基点移动路径的相反方向，也抵消。因此上下两条路径积分相同。与原分段光滑路径上的普通积分一致来自原函数端点公式。
\end{proof}
这个证明明确使用全纯性产生局部原函数，不要求同伦本身可微。对连续非全纯函数没有同样结论；\(\overline z\) 在不同圆周上的积分随面积改变。

\subsection{Cauchy 积分定理}
\label{b3:subsec:100202}

\begin{corollary}\label{b3:cor:10-null-homotopic-cauchy}
若闭路 \(\gamma\) 在 \(U\) 中可收缩，\(f\) 在 \(U\) 全纯，则 \(\int_\gamma f\,\mathrm{d}z=0\)。特别地，单连通区域中所有闭路积分为零。
\end{corollary}
\begin{proof}
将 \(\gamma\) 同伦到常路径，前一定理使积分等于常路径的零积分。
\end{proof}
可收缩性对闭路本身提出强条件。区域中有孔时，某些不能收缩的闭路仍可能对全部全纯函数积分为零，例如两个孔绕行的交换子圈。Cauchy 的同调形式将精确描述积分所需的较弱条件。

\ProseParagraphBreak
同伦积分不变性适用于同一全纯函数的整个同伦像。若把围道推过一个孤立奇点，函数在某个小格上没有原函数，内部边相消仍成立，但这个格的边界积分不再为0。留数定理正是把这项缺口计为 \(2\pi\mathrm{i}\) 乘留数。于是“可自由移动围道”应读成“在避开奇点的区域里移动”，跨过奇点的变化则需另加相应指数权重。

\ProseParagraphBreak
例如 \(f(z)=1/z\)，正向小圆与正向大圆在去心平面同伦，积分都为 \(2\pi\mathrm{i}\)；但不能把其中任何一个收缩成常路径。两者的差作为圈零同调，积分为0。这一例同时展示同伦积分不变、非零周期和差圈相消，三个结论并没有相互矛盾。源路径的端点、闭路的定向及函数的定义域共同决定哪条定理可用。

\subsection{原函数判别准则}
\label{b3:subsec:100203}

\begin{theorem}\label{b3:thm:10-primitive-periods}
区域 \(U\) 上全纯 \(f\) 有单值全纯原函数，当且仅当每条分段光滑闭路的积分为0。
\end{theorem}
\begin{proof}
有原函数时由端点公式得零周期。反向固定 \(a\)，令 \(F(z)=\int_a^z f\,\mathrm{d}\zeta\)。两路径拼成闭路，零周期使定义不依赖选择。在 \(z\) 附近小盘内，增加一段线段给 \(F(z+h)-F(z)=\int_{[z,z+h]}f\,\mathrm{d}\zeta\)，除以 \(h\) 并用连续性得 \(F'=f\)。
\end{proof}
周期是对闭路同伦类的加法映射：拼接相加，反向变号，所以它落在交换群 \((\mathbb C,+)\) 中。这已经预示非交换的基本群信息会被积分压缩。若只检查一个生成圈，应先证明所有圈由它生成，例如去心平面；在多孔域没有这种简化。

\ProseParagraphBreak
无零函数 \(g\) 有全纯对数当且仅当 \(g'/g\) 的全部周期为0。若只要求全纯 \(m\) 次根，则周期可以为 \(2\pi\mathrm{i}m\) 的整数倍，因为指数除以 \(m\) 后仍为1。对数要求周期本身为零，取根只要求整除，两项条件不能混淆。

\begin{proposition}[对数与根的周期条件]\label{b3:prop:10-log-root-periods}
设 \(g\) 在区域 \(U\) 全纯且无零。它有全纯对数，当且仅当 \(\int_\gamma g'/g\,\mathrm{d}z=0\) 对全部闭路成立。它有无零全纯 \(m\) 次根，当且仅当这些周期全属于 \(2\pi\mathrm{i}m\mathbb Z\)。
\end{proposition}
\begin{proof}
若 \(g=e^L\)，则 \(g'/g=L'\)，全部周期为0。反向用零周期取原函数 \(A\)，\((g e^{-A})'=0\)，故 \(g=c e^A\)、\(c\ne0\)，取一个常数对数即可。若 \(g=h^m\)，则 \(g'/g=m h'/h\)，而 \(\int h'/h\) 是 \(h\circ\gamma\) 对0的指数乘 \(2\pi\mathrm{i}\)。反向固定 \(a\)，选 \(c^m=g(a)\)，定义
\[
 h(z)=c\exp\left(\frac1m\int_a^z\frac{g'}g\,\mathrm{d}\zeta\right).
\]
两条路径的差是闭路，假设使指数变化为 \(2\pi\mathrm{i}\) 整数，所以 \(h\) 单值；局部原函数给它全纯无零。求导得 \((h^m/g)'=0\)，在 \(a\) 值为1，故 \(h^m=g\)。
\end{proof}
例如 \(g(z)=z^m\) 在圆环上有根 \(z\)，但没有对数；其基本周期为 \(2\pi\mathrm{i}m\)。\(g(z)=z\) 在圆环上无对数也没有 \(m>1\) 次根。无零条件与区域拓扑共同决定可否选分支，不能只检查表达式中的代数形式。

\ProseParagraphBreak
在有限连通的光滑平面域中，可选各内边界的正向小绕圈作为周期测试族。用不相交割线把各内边界连到外边界，切开后成单连通域；闭路经有限多次切线交叉分解，积分等于各基本周期的整数线性组合。这项构造要求边界分支和割线有限，且确实形成切开模型。一般有无限多孔的区域不能直接宣布“检查几个孔就够了”，应保留全部闭路条件或证明所选族生成了所需的同调。

\begin{proposition}[圆环周期与 Laurent 项]\label{b3:prop:10-annulus-period}
若 \(f\) 在圆环全纯，Laurent 展开为 \(\sum_{n\in\mathbb Z}a_nz^n\)，则任意闭路满足
\[
 \int_\gamma f\,\mathrm{d}z=2\pi\mathrm{i}a_{-1}\operatorname{Ind}(\gamma,0).
\]
特别地，\(f\) 在圆环有原函数当且仅当 \(a_{-1}=0\)。
\end{proposition}
\begin{proof}
闭路像紧，落在一个较小闭圆环中，Laurent 级数在其上绝对一致收敛，故逐项积分合法。每项 \(z^n\)、\(n\ne-1\) 有单值原函数 \(z^{n+1}/(n+1)\)，贡献0；\(z^{-1}\) 项由指数定义给所写式。若 \(a_{-1}=0\)，全部周期0，由原函数判别得存在；若有原函数，正向圆圈积分0使系数为0。
\end{proof}
这条结论并非说圆环里只有一个可能的极点。函数在圆环本身没有奇点，其 Laurent 负幂描述内孔一侧的解析行为，内孔中未必有一个单独可辨的奇点。周期只看 \(a_{-1}\)，其余无穷多个负幂在积分中都不见。这又展示积分比完整解析函数保留更少的信息。
